\textbf{（1）思路方面}

% 一
本项目首先针对海杂波强干扰下电磁频谱信号传播机理不清、微弱目标特征淹没于杂波背景的问题，在\textbf{“杂波解析微弱增强”}的指导思想下，探索复杂海洋环境下电磁频谱信号的时空耦合传播规律与干扰对消增强机理，为强杂波条件下微弱信号的精确提取提供坚实的理论支撑；
% 二
其次，针对舰载平台算力、存储与回传链路受限，海量频谱数据难以全量存储且可信难以保障问题，在\textbf{“端侧轻量可信存储”}的指导思想下，系统性开展事件驱动价值评估、表征压缩与链上链下协同存证一体化的轻量管理机理研究，为有限算力端侧频谱数据的高效存储与可信溯源提供方法基座；
% 三
最后，针对海域自主侦测中知识表征割裂、决策时效不足与多平台协同可信缺失的难题，在\textbf{“知识驱动可信决策”}的指导思想下，展开频谱知识库构建、边缘在线学习与分布式可信决策一体化框架研究，为无人平台集群的自主认知与协同决策提供“可表征、可演化、可信任”的全链条支撑。因此，本项目融合了电磁频谱感知、轻量数据管理与边缘智能决策方法，在研究思路上具有创新性。

\textbf{（2）技术方面}

% 在\textbf{“杂波解析微弱增强、端侧轻量可信存储、知识驱动可信决策”}的指导思想下，
通过深入探索和逐一突破，本项目有望取得以下三方面的技术创新：

% 一
\textbf{海杂波干扰下微弱电磁频谱信号传播规律分析与提取}：针对强海杂波背景下信号传播规律不清、微弱目标特征难以分离的难题，提出\textbf{电磁频谱信号传播规律、微弱目标频谱特征增强、电磁频谱信号提取方法}，突破强杂波与低信噪比下的特征淹没瓶颈，解决海上侦测“看不清、提不准”的难题，为高质量频谱特征输入提供机理支撑；

% 二
\textbf{面向有限算力端侧的频谱数据轻量存储与链上链下可信管理技术}：针对端侧资源受限下海量频谱数据存储开销大、可信溯源成本高的特征，提出\textbf{事件驱动的频谱数据价值评估方法、基于表征压缩的频谱向量数据库存储优化方法、链上链下协同的可追溯频谱数据可信管理方法}，突破高吞吐频谱数据的轻量存储与低成本存证瓶颈，解决端侧“存得下、查得快、信得过”的难题，为频谱数据全生命周期管理提供技术支撑；

% 三
\textbf{基于频谱知识驱动的边缘自主认知与分布式可信决策技术}：针对海域自主侦测中知识表征割裂、决策时效不足与多平台协同可信缺失的问题，提出\textbf{频谱知识库构建与动态管理方法、频谱知识驱动的自主决策与边缘在线学习方法、频谱认知协同与分布式可信决策方法}，突破\textbf{边缘算力约束下的认知演化与可信协同瓶颈}，解决无人平台集群“如何认知、如何决策、如何互信”的难题，为海域自主侦测任务的智能闭环提供一体化技术支撑。